\clearpage\pagestyle{otMain}\setlength{\otBlocEspacement}{2.5mm}\fontsize{11}{15}\selectfont\corrigesChapitre{19}{La proportionnalité}
\begin{blocLivre}{cor-19-1}\begin{corrige}{1}{Proportionnel ou pas ?}\demarche{Tester la phrase « fois plus » sur la situation.}
Les cahiers : $\res{\text{oui}}$, 2 fois plus de cahiers coûtent 2 fois plus. La taille : $\res{\text{non}}$, elle n’est pas multipliée par 2 quand l’âge double.
\end{corrige}\end{blocLivre}
\begin{blocLivre}{cor-19-2}\begin{corrige}{2}{Deux fois plus}\demarche{Doubler le prix quand la masse double.}
2 kg, c’est 2 fois plus : $2{,}40\times2=\res{4{,}80}$ euros. 4 kg, c’est encore le double : $\res{9{,}60}$ euros.
\end{corrige}\end{blocLivre}
\begin{blocLivre}{cor-19-3}\begin{corrige}{3}{La recette multipliée}\demarche{Comparer les nombres de personnes entre eux.}
8 personnes, 2 fois plus : $250\times2=\res{500}$ g. 12 personnes, 3 fois plus que 4 : $250\times3=\res{750}$ g.
\end{corrige}\end{blocLivre}
\begin{blocLivre}{cor-19-4}\begin{corrige}{4}{Les pains aux raisins}\demarche{Chercher le lien multiplicatif entre 4 et 12.}
12 pains, c’est $12:4=3$ fois plus de pains : ils coûtent 3 fois plus, soit $7\times3=\res{21}$ euros. Inutile de connaître le prix d’un pain.
\end{corrige}\end{blocLivre}
\begin{blocLivre}{cor-19-5}\begin{corrige}{5}{Les cahiers du bureau}\demarche{Multiplier le prix par le même nombre de fois.}
15 cahiers, c’est 3 fois plus : $8\times3=\res{24}$ euros.
\end{corrige}\end{blocLivre}
\begin{blocLivre}{cor-19-6}\begin{corrige}{6}{Fois moins}\demarche{Utiliser « fois plus » et « fois moins » depuis la même donnée.}
30 copies, 3 fois plus : $3\times3=\res9$ euros. 5 copies, 2 fois moins que 10 : $3:2=\res{1{,}50}$ euro.
\end{corrige}\end{blocLivre}
\begin{blocLivre}{cor-19-7}\begin{corrige}{7}{Le pot de peinture}\demarche{Relier chaque quantité de peinture à 6 litres.}
3 L, 2 fois moins : $30:2=\res{15}$ m². 18 L, 3 fois plus que 6 : $30\times3=\res{90}$ m².
\end{corrige}\end{blocLivre}
\begin{blocLivre}{cor-19-8}\begin{corrige}{8}{À allure régulière}\demarche{Comparer chaque durée à 2 heures.}
1 h, 2 fois moins de temps : $320:2=\res{160}$ km. 6 h, 3 fois plus que 2 h : $320\times3=\res{960}$ km. L’allure régulière rend la distance proportionnelle au temps.
\end{corrige}\end{blocLivre}
\begin{blocLivre}{cor-19-9}\begin{corrige}{9}{Les crêpes du goûter}\demarche{Raisonner dans le sens demandé par chaque question.}
500 g, 2 fois plus : $12\times2=\res{24}$ crêpes. 36 crêpes, 3 fois plus que 12 : $250\times3=\res{750}$ g de farine.
\end{corrige}\end{blocLivre}
\begin{blocLivre}{cor-19-10}\begin{corrige}{10}{L’abonnement piégé}\demarche{Calculer les deux montants avant de comparer.}
5 entrées : $10+5\times2=20$ euros. 10 entrées : $10+10\times2=30$ euros. Or $20\times2=40\ne30$ : Paul $\res{\text{se trompe}}$, la part fixe de l’abonnement casse la proportionnalité.
\end{corrige}\end{blocLivre}
\begin{blocLivre}{cor-19-11}\begin{corrige}{11}{La fontaine du square}\demarche{Faire le lien multiplicatif sur la grandeur connue.}
20 min, 4 fois plus : $40\times4=\res{160}$ L. Pour 120 L, soit 3 fois 40 L, il faut 3 fois plus de temps : $5\times3=\res{15}$ min.
\end{corrige}\end{blocLivre}
\begin{blocLivre}{cor-19-12}\begin{corrige}{12}{Le marché}\demarche{Relier chaque masse à 3 kg.}
9 kg, 3 fois plus : $7{,}50\times3=\res{22{,}50}$ euros. 6 kg, 2 fois plus : $7{,}50\times2=\res{15}$ euros.
\end{corrige}\end{blocLivre}
\begin{blocLivre}{cor-19-13}\begin{corrige}{13}{Deux boulangeries}\demarche{Ramener les deux offres au même nombre de croissants.}
Chez Bleuet, 8 croissants, 2 fois plus : $5\times2=10$ euros. Chez Coquelicot : $9{,}60$ euros. Comme $9{,}60<10$, c’est $\res{\text{chez Coquelicot}}$, avec $0{,}40$ euro d’écart.
\end{corrige}\end{blocLivre}
\begin{blocLivre}{cor-19-14}\begin{corrige}{14}{Le ruban au mètre}\demarche{Enchaîner un « fois moins » puis un « fois plus ».}
2 m, 2 fois moins que 4 m : $6:2=3$ euros. 10 m, 5 fois plus que 2 m : $3\times5=\res{15}$ euros. Passer par 2 m permet d’atteindre 10 m par multiplication.
\end{corrige}\end{blocLivre}
\begin{blocLivre}{cor-19-15}\begin{corrige}{15}{Le carré qui grandit}\demarche{Tester chaque grandeur séparément.}
Périmètre : $4\times3=12$ devient $4\times6=24$ : il $\res{\text{double}}$, le périmètre est proportionnel au côté. Aire : $3\times3=9$ carreaux devient $6\times6=36$ : elle est multipliée par $\res4$, pas par 2 : l’aire n’est pas proportionnelle au côté.
\end{corrige}\end{blocLivre}
